% ============================================================
% MÓDULO: FINANCIAMIENTO DEL EMPRENDIMIENTO Y PROCESO DE COMPRAS
% ============================================================
\clearpage
\section*{MÓDULO: FINANCIAMIENTO Y PROCESO DE COMPRAS}
\subsection*{De dónde sale el dinero y cómo se compra en una empresa}

\begin{recuadro}
\textbf{Objetivo del módulo:} conocer las distintas fuentes de financiamiento de un
emprendimiento (además del préstamo bancario) y el circuito formal de compras de una empresa,
desde el pedido de presupuesto hasta el pago de la factura.
\end{recuadro}

% ------------------------------------------------------------
\section*{1. FUENTES DE FINANCIAMIENTO}

\begin{center}
\small
\begin{tabular}{|p{0.22\linewidth}|p{0.42\linewidth}|p{0.26\linewidth}|}
\hline
\textbf{Fuente} & \textbf{Cómo funciona} & \textbf{Para qué etapa} \\ \hline
Ahorros propios (\textit{bootstrapping}) & Se arranca con recursos propios y se reinvierten
las ganancias. Se conserva el control total. & Inicio \\ \hline
Familia y amigos & Préstamos o aportes de personas cercanas. Conviene dejarlo por escrito. &
Inicio \\ \hline
Clientes & Anticipos o señas por los trabajos; preventa. & Cualquiera \\ \hline
Proveedores & Plazo para pagar las compras (30, 60 días). & Cualquiera \\ \hline
Préstamos bancarios & Personales, para pymes, prendarios, con tasa y cuotas (ver unidad de
préstamos). & Con ventas \\ \hline
SGR (Sociedades de Garantía Recíproca) & Dan un aval a la pyme para que consiga crédito en
mejores condiciones. & Pyme en marcha \\ \hline
Leasing & Se alquila un equipo (por ejemplo, un banco de pruebas) con opción de compra al
final. & Inversión en equipos \\ \hline
Programas públicos, subsidios y concursos & Aportes no reembolsables o créditos blandos de
nación, provincia o municipio; premios a emprendedores. & Inicio y crecimiento \\ \hline
Inversores ángeles y capital de riesgo & Invierten a cambio de una participación en la
empresa; buscan alto crecimiento. & Proyectos escalables \\ \hline
Financiamiento colectivo (crowdfunding) & Muchas personas aportan pequeños montos a través de
una plataforma. & Inicio y lanzamiento \\ \hline
\end{tabular}
\end{center}

\subsubsection*{1.1. El crowdfunding}
El \textbf{financiamiento colectivo} reúne pequeños aportes de muchas personas por internet.
Hay distintos tipos:
\begin{itemize}
  \item \textbf{De recompensa:} los aportantes reciben el producto (preventa) o un beneficio.
        Útil para validar si hay interés antes de fabricar.
  \item \textbf{De donación:} para proyectos sociales o culturales.
  \item \textbf{De inversión (\textit{equity crowdfunding}):} los aportantes reciben
        acciones o títulos de la empresa. En Argentina lo regula la \textbf{Ley 27.349 de
        Apoyo al Capital Emprendedor}, se hace a través de \textbf{Plataformas de
        Financiamiento Colectivo (PFC)} autorizadas y controladas por la \textbf{Comisión
        Nacional de Valores (CNV)}. En 2026 la CNV simplificó el régimen (RG 1125/2026).
\end{itemize}

\paragraph*{Cómo elegir.}
Cada fuente tiene un costo: los préstamos cobran \textbf{intereses}; los inversores se quedan
con una \textbf{parte de la empresa} y de las decisiones; los subsidios exigen
\textbf{rendir} cómo se usó el dinero. Regla práctica: no financiar con deuda de corto plazo
una inversión que tarda años en recuperarse, y no endeudarse en dólares si se cobra en pesos.

% ------------------------------------------------------------
\section*{2. EL PROCESO DE COMPRAS}

En una empresa no se compra ``yendo al negocio'': hay un circuito con controles para evitar
errores, fraudes y compras innecesarias, y para que cada gasto quede justificado.

\begin{enumerate}
  \item \textbf{Solicitud de compra (pedido interno):} el área que necesita algo (el taller, el
        pañol) pide el material indicando código, descripción, cantidad, fecha requerida y
        orden de trabajo o centro de costo. Muchas veces se genera automáticamente cuando el
        stock baja del mínimo.
  \item \textbf{Solicitud de presupuesto (cotización):} compras pide precio a varios
        proveedores (al menos tres, si es posible) con las mismas especificaciones.
  \item \textbf{Cuadro comparativo:} se comparan precio, plazo de entrega, forma de pago,
        garantía, calidad y trazabilidad. No siempre gana el más barato.
  \item \textbf{Orden de compra (OC):} documento formal que la empresa envía al proveedor
        elegido. Indica proveedor, artículos, cantidades, precios, condiciones de pago, lugar y
        fecha de entrega. Es un compromiso: una vez aceptada, obliga a ambas partes.
  \item \textbf{Recepción:} al llegar la mercadería con el \textbf{remito}, se controla
        cantidad, estado y documentación, y se firma la recepción. En el ERP se valida la
        entrada al stock.
  \item \textbf{Factura del proveedor:} se recibe y se registra.
  \item \textbf{Control de tres vías:} antes de pagar se verifica que coincidan la
        \textbf{orden de compra}, el \textbf{remito/recepción} y la \textbf{factura}
        (cantidades y precios).
  \item \textbf{Pago:} según las condiciones pactadas (contado, 30 días), con su
        \textbf{orden de pago} y el recibo del proveedor. Se aplican las retenciones
        impositivas que correspondan.
\end{enumerate}

\begin{recuadro}
\textbf{Compras en aviación.} Un componente aeronáutico no se puede comprar a cualquier
proveedor ni solo por precio. Debe venir de un \textbf{proveedor aprobado} y con su
\textbf{documentación de aeronavegabilidad} (por ejemplo, el formulario FAA 8130-3, el EASA
Form 1 o el certificado de conformidad del fabricante), que garantiza su origen y estado. Sin
esa documentación, la pieza no puede instalarse en una aeronave. El área de compras y el pañol
verifican estos papeles en la recepción y los archivan para mantener la
\textbf{trazabilidad}.
\end{recuadro}

% ------------------------------------------------------------
\section*{3. ACTIVIDADES DEL MÓDULO}

\begin{enumerate}
  \item \textbf{Plan de financiamiento:} para el proyecto final, decidir qué parte de la
        inversión se cubre con cada fuente y justificar.
  \item \textbf{Investigación:} buscar un programa vigente de apoyo a emprendedores (nacional,
        provincial o municipal) y un caso de crowdfunding exitoso; presentar requisitos y
        resultados.
  \item \textbf{Circuito de compras completo:} a partir de un pedido del taller, solicitar tres
        presupuestos (reales o simulados), armar el cuadro comparativo, emitir la orden de
        compra, registrar la recepción y controlar la factura. Hacerlo también en Odoo dentro
        del proyecto integrador.
  \item \textbf{Detectar errores:} con un juego de OC, remito y factura que no coinciden,
        identificar las diferencias y redactar el reclamo al proveedor.
\end{enumerate}
